\documentclass[12pt,a4paper]{article}
\usepackage[utf8]{inputenc}
\usepackage[spanish,german]{babel}
\usepackage[T1]{fontenc}
\usepackage{geometry}
\geometry{a4paper, margin=2.5cm}
\usepackage{titlesec}
\usepackage{enumitem}
\usepackage{multicol}
\usepackage{csquotes}

\titleformat{\section}
{\Large\bfseries}
{}
{0em}
{}[\titlerule]

\titleformat{\subsection}
{\large\bfseries}
{}
{0em}
{}

\setlength{\parindent}{0pt}
\setlength{\parskip}{0.5em}

\begin{document}

\title{\textbf{DIALOG: Hablar de Medio Ambiente}}
\author{Spanisch A1/A2}
\date{}
\maketitle

\textbf{Tema:} Hablar de medio ambiente \\
\textbf{Nivel:} A2 \\
\textbf{Complejidad:} Mittel \\
\textbf{Palabras:} 842

%\vspace{1em}

\textit{Nota: Se recomienda revisar primero el vocabulario antes de leer el diálogo.}

%\vspace{1em}

\section*{Diálogo}

\textbf{Ana:} Hola, Diego. ¿Cómo estás?

\textbf{Diego:} Hola, Ana. Estoy bien, gracias. ¿Y tú?

\textbf{Ana:} Muy bien. Oye, ¿sabes qué? Ayer fui al parque y vi mucha basura por todas partes. Me da mucha tristeza.

\textbf{Diego:} Sí, es un problema muy serio. Por eso yo siempre llevo una bolsa para recoger basura cuando salgo a caminar.

\textbf{Ana:} ¡Qué buena idea! Yo también quiero hacer algo por el medio ambiente. ¿Qué más podemos hacer para ayudar?

\textbf{Diego:} Bueno, hay muchas cosas. Por ejemplo, podemos usar menos plástico. Yo ahora uso una botella de agua reutilizable para no comprar botellas de plástico.

\textbf{Ana:} Eso es excelente. Yo también compré una botella reutilizable la semana pasada. La uso todos los días para ir al trabajo.

\textbf{Diego:} Perfecto. Otra cosa importante es reciclar. En mi casa separamos la basura por tipo: papel, plástico y vidrio.

\textbf{Ana:} Nosotros también reciclamos. Mi mamá compró contenedores especiales para cada tipo de basura. Los pusimos en la cocina.

\textbf{Diego:} Muy bien. ¿Y qué haces para ahorrar energía en tu casa?

\textbf{Ana:} Bueno, apagamos las luces cuando no las necesitamos. También usamos bombillas de bajo consumo. ¿Y tú?

\textbf{Diego:} Yo hago lo mismo. Además, desconecto los aparatos electrónicos por la noche para ahorrar energía.

\textbf{Ana:} Eso es muy inteligente. Oye, ¿sabes que el próximo sábado hay una limpieza del río? Están buscando voluntarios para ayudar.

\textbf{Diego:} ¡Qué interesante! ¿A qué hora es? Me gustaría participar.

\textbf{Ana:} Es a las nueve de la mañana. Van a estar allí por tres horas, hasta el mediodía. ¿Quieres venir conmigo?

\textbf{Diego:} ¡Claro que sí! Es una buena oportunidad para hacer algo por nuestra comunidad. ¿Dónde nos encontramos?

\textbf{Ana:} Podemos encontrarnos en la entrada del parque a las ocho y media. Así llegamos juntos al río.

\textbf{Diego:} Perfecto. Voy a llevar guantes y bolsas para recoger la basura.

\textbf{Ana:} Excelente. Yo también voy a llevar todo lo necesario. Por cierto, ¿has visto el documental sobre el cambio climático que pasaron en la televisión?

\textbf{Diego:} No, no lo vi. ¿De qué se trata?

\textbf{Ana:} Habla sobre los problemas del medio ambiente y qué podemos hacer para solucionarlos. Es muy educativo. Lo pasan por la televisión todos los domingos por la noche.

\textbf{Diego:} Voy a verlo este domingo entonces. Me interesa mucho el tema. ¿Sabes? Creo que todos tenemos que hacer algo para proteger el planeta.

\textbf{Ana:} Tienes toda la razón. Cada pequeña acción cuenta. Por ejemplo, caminar o usar la bicicleta en lugar del auto ayuda mucho.

\textbf{Diego:} Sí, yo ahora voy al trabajo en bicicleta tres veces por semana. Es bueno para el medio ambiente y también para mi salud.

\textbf{Ana:} ¡Qué bien! Yo todavía no tengo bicicleta, pero camino mucho. Voy a pie al supermercado y a otros lugares cercanos.

\textbf{Diego:} Eso también ayuda. Otra cosa que podemos hacer es comprar productos locales. Así no se necesita transportar la comida desde muy lejos.

\textbf{Ana:} Sí, es verdad. Yo compro frutas y verduras en el mercado local. Son más frescas y ayudan a los productores de la zona.

\textbf{Diego:} Exacto. Y también podemos reducir el consumo de carne. No digo que todos tengan que ser vegetarianos, pero comer menos carne es bueno para el medio ambiente.

\textbf{Ana:} Tienes razón. Yo como carne solo dos o tres veces por semana ahora. El resto del tiempo como más vegetales y legumbres.

\textbf{Diego:} Eso es perfecto. Mira, creo que lo más importante es que todos hagamos pequeños cambios en nuestra vida diaria. No necesitamos hacer todo perfecto, pero cada acción cuenta.

\textbf{Ana:} Estoy completamente de acuerdo. Y es bueno hablar sobre esto con otras personas para crear conciencia.

\textbf{Diego:} Sí, por eso me gusta participar en actividades como la limpieza del río. Es una forma de hacer algo concreto y también de conocer a otras personas que se preocupan por el medio ambiente.

\textbf{Ana:} Me parece genial. Bueno, entonces nos vemos el sábado a las ocho y media en la entrada del parque.

\textbf{Diego:} Perfecto. Nos vemos el sábado. ¡Que tengas un buen día!

\textbf{Ana:} Gracias, tú también. ¡Hasta el sábado!

\newpage

\section*{Vocabulario}

\begin{multicols}{2}
\begin{itemize}[leftmargin=*]
    \item \textbf{medio ambiente} - Umwelt
    \item \textbf{basura} - Müll, Abfall
    \item \textbf{por todas partes} - überall
    \item \textbf{recoger} - sammeln, aufheben
    \item \textbf{reutilizable} - wiederverwendbar
    \item \textbf{reciclar} - recyceln
    \item \textbf{separar} - trennen, sortieren
    \item \textbf{contenedor} - Behälter, Container
    \item \textbf{ahorrar} - sparen
    \item \textbf{energía} - Energie
    \item \textbf{bombilla} - Glühbirne
    \item \textbf{bajo consumo} - niedriger Verbrauch
    \item \textbf{desconectar} - abstecken, trennen
    \item \textbf{aparato electrónico} - elektronisches Gerät
    \item \textbf{voluntario} - Freiwilliger/Freiwillige
    \item \textbf{participar} - teilnehmen
    \item \textbf{oportunidad} - Gelegenheit, Chance
    \item \textbf{comunidad} - Gemeinschaft
    \item \textbf{guantes} - Handschuhe
    \item \textbf{cambio climático} - Klimawandel
    \item \textbf{educativo} - lehrreich, bildend
    \item \textbf{proteger} - schützen
    \item \textbf{planeta} - Planet
    \item \textbf{producto local} - lokales Produkt
    \item \textbf{transportar} - transportieren
    \item \textbf{productor} - Produzent, Erzeuger
    \item \textbf{reducir} - reduzieren, verringern
    \item \textbf{consumo} - Verbrauch, Konsum
    \item \textbf{vegetariano} - Vegetarier
    \item \textbf{legumbre} - Hülsenfrucht
    \item \textbf{conciencia} - Bewusstsein
    \item \textbf{concreto} - konkret
\end{itemize}
\end{multicols}

\newpage

\section*{Preguntas de Comprensión}

\begin{enumerate}[leftmargin=*]
    \item ¿Qué vio Ana en el parque que le dio tristeza?
    \item ¿Qué lleva Diego cuando sale a caminar?
    \item ¿Qué tipo de botella usa Diego para el agua?
    \item ¿Cómo separan la basura en la casa de Diego?
    \item ¿Qué hace Ana para ahorrar energía en su casa?
    \item ¿Qué hace Diego con los aparatos electrónicos por la noche?
    \item ¿Cuándo es la limpieza del río?
    \item ¿Dónde se van a encontrar Ana y Diego?
    \item ¿Qué va a llevar Diego a la limpieza del río?
    \item ¿Cuándo pasan el documental sobre cambio climático en la televisión?
    \item ¿Cuántas veces por semana va Diego al trabajo en bicicleta?
    \item ¿Dónde compra Ana frutas y verduras?
    \item ¿Cuántas veces por semana come Ana carne ahora?
    \item ¿Por qué le gusta a Diego participar en actividades como la limpieza del río?
    \item ¿Qué hora es la limpieza del río?
    \item ¿Qué compró la mamá de Ana para reciclar?
    \item ¿Qué usa Ana en lugar de la bicicleta para moverse?
    \item ¿Por qué es bueno comprar productos locales según Diego?
    \item ¿Qué opina Diego sobre comer carne?
    \item ¿Qué es lo más importante según Diego para ayudar al medio ambiente?
\end{enumerate}

\newpage

\section*{Verbo Irregular: HACER (machen, tun)}

\textit{Nota: \enquote{vosotros/as} se usa solo en España, no en español sudamericano. Se incluye aquí para completitud.}

\begin{multicols}{2}

\subsection*{Presente}
\begin{itemize}[leftmargin=*]
    \item yo - hago
    \item tú - haces
    \item él/ella/usted - hace
    \item nosotros/as - hacemos
    \item vosotros/as - hacéis
    \item ustedes - hacen
    \item ellos/ellas - hacen
\end{itemize}

\subsection*{Pretérito Indefinido}
\begin{itemize}[leftmargin=*]
    \item yo - hice
    \item tú - hiciste
    \item él/ella/usted - hizo
    \item nosotros/as - hicimos
    \item vosotros/as - hicisteis
    \item ustedes - hicieron
    \item ellos/ellas - hicieron
\end{itemize}

\subsection*{Pretérito Imperfecto}
\begin{itemize}[leftmargin=*]
    \item yo - hacía
    \item tú - hacías
    \item él/ella/usted - hacía
    \item nosotros/as - hacíamos
    \item vosotros/as - hacíais
    \item ustedes - hacían
    \item ellos/ellas - hacían
\end{itemize}

\columnbreak

\subsection*{Futuro Simple}
\begin{itemize}[leftmargin=*]
    \item yo - haré
    \item tú - harás
    \item él/ella/usted - hará
    \item nosotros/as - haremos
    \item vosotros/as - haréis
    \item ustedes - harán
    \item ellos/ellas - harán
\end{itemize}

\subsection*{Pretérito Perfecto}
\begin{itemize}[leftmargin=*]
    \item yo - he hecho
    \item tú - has hecho
    \item él/ella/usted - ha hecho
    \item nosotros/as - hemos hecho
    \item vosotros/as - habéis hecho
    \item ustedes - han hecho
    \item ellos/ellas - han hecho
\end{itemize}

\end{multicols}

\newpage

\section*{Explicación Gramatical}

\textbf{Por vs. Para -- Der Unterschied zwischen \enquote{por} und \enquote{para}}

Die Präpositionen \enquote{por} und \enquote{para} sind im Spanischen oft verwirrend, weil beide im Deutschen oft mit \enquote{für} übersetzt werden können. Es gibt jedoch wichtige Unterschiede:

\textbf{\enquote{POR} wird verwendet für:}
\begin{itemize}[leftmargin=*]
    \item \textbf{Grund/Ursache}: \enquote{por eso} (deshalb), \enquote{por la noche} (in der Nacht)
    \item \textbf{Zeitdauer}: \enquote{por tres horas} (drei Stunden lang)
    \item \textbf{Art und Weise}: \enquote{por tipo} (nach Art/Typ)
    \item \textbf{Ersatz/Vertretung}: \enquote{por nuestra comunidad} (für unsere Gemeinschaft)
    \item \textbf{Bewegung durch einen Ort}: \enquote{por todas partes} (überall)
    \item \textbf{Preis}: \enquote{por la mañana} (am Morgen -- Zeitpunkt)
\end{itemize}

\textbf{\enquote{PARA} wird verwendet für:}
\begin{itemize}[leftmargin=*]
    \item \textbf{Zweck/Ziel}: \enquote{para ayudar} (um zu helfen), \enquote{para recoger basura} (um Müll zu sammeln)
    \item \textbf{Bestimmung/Zielort}: \enquote{para el trabajo} (für die Arbeit), \enquote{para ir al trabajo} (um zur Arbeit zu gehen)
    \item \textbf{Zeitpunkt in der Zukunft}: \enquote{para el sábado} (für Samstag)
    \item \textbf{Meinung}: \enquote{para mí} (für mich)
    \item \textbf{Vergleich}: \enquote{para un estudiante} (für einen Studenten)
\end{itemize}

\textbf{Ejemplos del texto:}
\begin{itemize}[leftmargin=*]
    \item \enquote{vi mucha basura \textbf{por} todas partes} -- Bewegung durch einen Raum
    \item \enquote{\textbf{Por} eso yo siempre llevo una bolsa} -- Grund/Ursache
    \item \enquote{una botella de agua reutilizable \textbf{para} no comprar botellas} -- Zweck
    \item \enquote{separamos la basura \textbf{por} tipo} -- Art und Weise
    \item \enquote{contenedores especiales \textbf{para} cada tipo} -- Zweck/Bestimmung
    \item \enquote{desconecto los aparatos \textbf{por} la noche} -- Zeitpunkt
    \item \enquote{están allí \textbf{por} tres horas} -- Zeitdauer
    \item \enquote{hacer algo \textbf{por} nuestra comunidad} -- für/zu Gunsten von
    \item \enquote{guantes y bolsas \textbf{para} recoger la basura} -- Zweck
    \item \enquote{lo pasan \textbf{por} la televisión} -- Medium/Kanal
    \item \enquote{todos los domingos \textbf{por} la noche} -- Zeitpunkt
    \item \enquote{hacer algo \textbf{para} proteger el planeta} -- Zweck
    \item \enquote{en bicicleta \textbf{para} ir al trabajo} -- Zweck
    \item \enquote{comprar productos locales \textbf{para} ayudar} -- Zweck
\end{itemize}

\textbf{Tipp}: Wenn du \enquote{um zu} sagen kannst, dann ist es meistens \enquote{para}. Wenn es um Grund, Zeitdauer oder Bewegung geht, dann ist es \enquote{por}.

\end{document}
